\documentclass[11pt,a4paper]{article}
% preamble.tex — shared setup for human-verification documents
% Compile each document with:  tectonic <File>.tex

\usepackage[T1]{fontenc}
\usepackage[utf8]{inputenc}
\usepackage{lmodern}
\usepackage{amssymb,amsmath,amsthm}
\usepackage{mathtools}
\usepackage[dvipsnames]{xcolor}
\usepackage{listings}
\usepackage{hyperref}
\usepackage[margin=2.5cm]{geometry}
\usepackage{parskip}

% ── Colours ──────────────────────────────────────────────────────────
\definecolor{leanBg}{HTML}{F7F7F0}
\definecolor{haskellBg}{HTML}{F0F4F7}
\definecolor{keyword}{HTML}{7B3294}
\definecolor{comment}{HTML}{5A8C32}
\definecolor{noteFg}{HTML}{B35900}

% ── Lean 4 listings language ─────────────────────────────────────────
% NOTE: Lean listing blocks use ASCII approximations of unicode symbols
% (e.g. -> for arrow, o. for compose).  The literate key substitutes
% these with proper math-mode renderings in the PDF.
\lstdefinelanguage{Lean4}{
  morekeywords={def,theorem,structure,where,namespace,end,import,fun,
                let,have,show,by,exact,intro,cases,rfl,subst,obtain,
                funext,simp,rw,apply,variable,Type,Prop,instance,
                section,open,class,extends,deriving,sorry,noncomputable},
  sensitive=true,
  morecomment=[l]{--},
  morecomment=[n]{/-}{-/},
  morestring=[b]",
}

% ── Lean listing style ───────────────────────────────────────────────
\lstdefinestyle{lean}{
  language=Lean4,
  backgroundcolor=\color{leanBg},
  basicstyle=\small\ttfamily,
  keywordstyle=\color{keyword}\bfseries,
  commentstyle=\color{comment}\itshape,
  breaklines=true,
  frame=single,
  framerule=0.4pt,
  xleftmargin=4pt,
  xrightmargin=4pt,
  aboveskip=6pt,
  belowskip=4pt,
  literate=
    {->}{{$\to$}}2
    {=>}{{$\Rightarrow$}}2
    {<|}{{$\langle$}}1
    {|>}{{$\rangle$}}1
    {==>}{{$\Longrightarrow$}}3
    { o. }{{$\;\circ\;$}}3
    {forall}{{$\forall$}}1
    {:=}{{$\,:=\,$}}2
    {!=}{{$\neq$}}2
    {<=}{{$\leq$}}2
    {>=}{{$\geq$}}2,
}

% ── Haskell listing style ───────────────────────────────────────────
\lstdefinestyle{haskell}{
  language=Haskell,
  backgroundcolor=\color{haskellBg},
  basicstyle=\small\ttfamily,
  keywordstyle=\color{keyword}\bfseries,
  commentstyle=\color{comment}\itshape,
  breaklines=true,
  frame=single,
  framerule=0.4pt,
  xleftmargin=4pt,
  xrightmargin=4pt,
  aboveskip=6pt,
  belowskip=4pt,
  literate=
    {->}{{$\to$}}2
    {=>}{{$\Rightarrow$}}2,
}

% ── Theorem environments ────────────────────────────────────────────
\theoremstyle{definition}
\newtheorem{defn}{Definition}[section]
\newtheorem{thm}[defn]{Theorem}
\newtheorem{prop}[defn]{Proposition}

\theoremstyle{remark}
\newtheorem*{fidelitynote}{Fidelity note}

% ── Convenience macros ───────────────────────────────────────────────
\newcommand{\lean}[1]{\texttt{#1}}      % inline Lean identifier
\newcommand{\hask}[1]{\texttt{#1}}      % inline Haskell identifier
\newcommand{\Ext}{\mathrm{Ext}}
\newcommand{\Pos}{\mathrm{Pos}}
\newcommand{\Shape}{\mathrm{Shape}}


% Double-bracket notation for the extension functor
\usepackage{stmaryrd}

\title{Human Verification: \texttt{Containers.FourMonoidal}\\[4pt]
  \large The Four Monoidal Structures on $\mathbf{Cont}$}
\author{Machine-generated from \texttt{lean/FourMonoidal.lean} for review by Neil Ghani}
\date{July 2026}

\begin{document}
\maketitle

\begin{abstract}
This document presents every definition and theorem \emph{statement} from the
Lean~4 module \lean{Containers.FourMonoidal} in three parallel forms: the original
Lean, a readable Haskell rendering, and precise mathematical English.  Proofs
are omitted --- they are machine-checked by Lean.  The purpose is to let a
human verifier confirm that the formalisation captures the intended
mathematics.

The file completes the account of all four monoidal structures on $\mathbf{Cont}$:
the coproduct $(+, 0)$, the categorical product $(\times, 1)$, the Dirichlet
tensor $(\otimes, y)$ (in \texttt{Dirichlet.lean}), and the sequential operator
$(\triangleleft, I)$ (in \texttt{Monoidal.lean}).  It also proves that the
extension functor $\llbracket{-}\rrbracket$ is \emph{strong} monoidal for $+$
and $\times$, and provides the comparison isomorphisms.
\end{abstract}

\tableofcontents
\bigskip

%===================================================================
\section{Overview}
%===================================================================

The four monoidal structures on $\mathbf{Cont}$ are summarised by the table:

\begin{center}
\begin{tabular}{llll}
\hline
Structure & Shapes of $C \otimes D$ & Positions of $C \otimes D$ & Unit \\
\hline
coproduct $+$ & $S \sqcup T$ & $[P, Q]$ (by injection) & $0 = \emptyset \triangleleft \cdots$ \\
product $\times$ & $S \times T$ & $P\,s \sqcup Q\,t$ & $1 = \mathbf{1} \triangleleft \emptyset$ \\
Dirichlet $\otimes$ & $S \times T$ & $P\,s \times Q\,t$ & $y = \mathbf{1} \triangleleft \mathbf{1}$ \\
sequential $\triangleleft$ & $\Sigma\,s, P\,s \to T$ & $\Sigma\,p, Q(f(p))$ & $I = y$ \\
\hline
\end{tabular}
\end{center}

\medskip
\noindent
\textbf{Key observation.} The product $\times$ and Dirichlet $\otimes$ have the
same shapes $S \times T$ but differ only in their positions (sum vs.\ product).
Yet the \emph{extensions} invert the reading:
$\llbracket C \times D \rrbracket X \cong \llbracket C \rrbracket X \times \llbracket D \rrbracket X$
(pointwise product), while
$\llbracket C \otimes D \rrbracket X$ is \emph{not} any pointwise combination
of $\llbracket C \rrbracket X$ and $\llbracket D \rrbracket X$.
This is the exponential law: $y^{A+B} \cong y^A \times y^B$, but $y^{A \times B}$
admits no such splitting.

Only $\lean{ContMonoidal}$ (sequential $\triangleleft$) is the canonical
\lean{MonoidalCategory Container} \emph{instance}; the other three are named
\lean{def}s to avoid instance clashes.

%===================================================================
\section{Unit containers}
%===================================================================

\subsection{Definition: \lean{Container.zero}}

\begin{defn}[The empty container $0$]
The \textbf{empty container} $0 = \emptyset \triangleleft (\mathrm{elim})$ has
no shapes and (vacuously) no positions.  It is the unit for the coproduct $+$
and the initial object of $\mathbf{Cont}$.
Its extension $\llbracket 0 \rrbracket X$ is the empty type for every $X$.
\end{defn}

\paragraph{Lean.}
\begin{lstlisting}[style=lean]
def Container.zero : Container where
  Shape := Empty
  Pos   := fun e => e.elim
\end{lstlisting}

\paragraph{Haskell.}
\begin{lstlisting}[style=haskell]
zeroContainer :: Container Void p
zeroContainer = Container { shape = absurd, pos = absurd }
-- Extension: Ext ZeroContainer x  is  uninhabited.
\end{lstlisting}


\subsection{Definition: \lean{Container.one}}

\begin{defn}[The one-shape container $1$]
The \textbf{one-shape, no-position container} $1 = \mathbf{1} \triangleleft \emptyset$
is the polynomial $y^0$.  It is the unit for the categorical product $\times$
and the terminal object of $\mathbf{Cont}$.
Its extension $\llbracket 1 \rrbracket X$ is a singleton type for every $X$.
\end{defn}

\paragraph{Lean.}
\begin{lstlisting}[style=lean]
def Container.one : Container where
  Shape := Unit
  Pos   := fun _ => Empty
\end{lstlisting}

\paragraph{Haskell.}
\begin{lstlisting}[style=haskell]
oneContainer :: Container () Void
oneContainer = Container { shape = const (), pos = const absurd }
-- Extension: Ext OneContainer x  =~=  Unit.
\end{lstlisting}


%===================================================================
\section{The coproduct \texorpdfstring{$(+, 0)$}{(+,0)}}
%===================================================================

\subsection{Definition: \lean{Container.coprod\textsubscript{2}}}

\begin{defn}[Bifunctorial action of $+$]
The morphism action $\lean{coprod}_2\,\mu\,\nu : C + D \Rightarrow C' + D'$
is derived from the coproduct's universal property: apply $\mu$ to $C$-shapes
and $\nu$ to $D$-shapes via the injections, then co-pair.
\end{defn}

\paragraph{Lean.}
\begin{lstlisting}[style=lean]
def Container.coprod2 {C C' D D' : Container}
    (mu : ContainerMorphism C C') (nu : ContainerMorphism D D') :
    ContainerMorphism (C.coprod D) (C'.coprod D') :=
  (mu.comp (Container.inl C' D')).desc (nu.comp (Container.inr C' D'))
\end{lstlisting}

\paragraph{Haskell.}
\begin{lstlisting}[style=haskell]
coprod2 :: Morph c p c' p' -> Morph d q d' q'
        -> Morph (c `Coprod` d) (Either p q) (c' `Coprod` d') (Either p' q')
coprod2 mu nu = desc (mu `compM` inlMorph) (nu `compM` inrMorph)
\end{lstlisting}


\subsection{Theorem: \lean{Container.coprod\textsubscript{2}\_id}}

\begin{thm}[Coproduct bifunctoriality: identity]
$\lean{coprod}_2\,(\mathrm{id}_C)\,(\mathrm{id}_D) = \mathrm{id}_{C+D}$.
\end{thm}

\paragraph{Lean.}
\begin{lstlisting}[style=lean]
theorem Container.coprod2_id (C D : Container) :
    Container.coprod2 (ContainerMorphism.id C) (ContainerMorphism.id D)
      = ContainerMorphism.id (C.coprod D)
\end{lstlisting}

\paragraph{Haskell.}
\begin{lstlisting}[style=haskell]
-- coprod2 idMorph idMorph = idMorph
\end{lstlisting}


\subsection{Theorem: \lean{Container.coprod\textsubscript{2}\_comp}}

\begin{thm}[Coproduct bifunctoriality: interchange]
$\lean{coprod}_2\,(\mu \cdot \mu')\,(\nu \cdot \nu')
= \lean{coprod}_2\,\mu\,\nu \;\cdot\; \lean{coprod}_2\,\mu'\,\nu'$.
\end{thm}

\paragraph{Lean.}
\begin{lstlisting}[style=lean]
theorem Container.coprod2_comp {C C' C'' D D' D'' : Container}
    (mu : ContainerMorphism C C') (mu' : ContainerMorphism C' C'')
    (nu : ContainerMorphism D D') (nu' : ContainerMorphism D' D'') :
    Container.coprod2 (mu.comp mu') (nu.comp nu')
      = (Container.coprod2 mu nu).comp (Container.coprod2 mu' nu')
\end{lstlisting}

\paragraph{Haskell.}
\begin{lstlisting}[style=haskell]
-- coprod2 (mu `compM` mu') (nu `compM` nu')
--   = coprod2 mu nu `compM` coprod2 mu' nu'
\end{lstlisting}


\subsection{Definition: \lean{Container.coprodAssociator}}

\begin{defn}[Coproduct associator]
The isomorphism $(C + D) + E \cong C + (D + E)$:
\begin{itemize}
  \item \emph{Shapes}: $\mathrm{Sum}$-reassociation
    $(\mathrm{inl}\,(\mathrm{inl}\,a) \mapsto \mathrm{inl}\,a$, etc.).
  \item \emph{Positions}: identity in every branch.
\end{itemize}
\end{defn}

\paragraph{Lean.}
\begin{lstlisting}[style=lean]
def Container.coprodAssociator (C D E : Container) :
    ContainerIso ((C.coprod D).coprod E) (C.coprod (D.coprod E)) where
  hom :=
    { onShapes := fun x => match x with
        | Sum.inl (Sum.inl a) => Sum.inl a
        | Sum.inl (Sum.inr b) => Sum.inr (Sum.inl b)
        | Sum.inr c           => Sum.inr (Sum.inr c)
      onPos := fun x => match x with
        | Sum.inl (Sum.inl _) => id
        | Sum.inl (Sum.inr _) => id
        | Sum.inr _           => id }
  inv :=
    { onShapes := fun x => match x with
        | Sum.inl a           => Sum.inl (Sum.inl a)
        | Sum.inr (Sum.inl b) => Sum.inl (Sum.inr b)
        | Sum.inr (Sum.inr c) => Sum.inr c
      onPos := fun x => match x with
        | Sum.inl _           => id
        | Sum.inr (Sum.inl _) => id
        | Sum.inr (Sum.inr _) => id }
  hom_inv := ...   -- proof omitted
  inv_hom := ...
\end{lstlisting}

\paragraph{Haskell.}
\begin{lstlisting}[style=haskell]
coprodAssociator :: ContainerIso
  (Coprod (Coprod c d) e) (Coprod c (Coprod d e))
coprodAssociator = ...
-- Shapes: standard Either reassociation; positions: id everywhere.
\end{lstlisting}


\subsection{Definition: \lean{Container.coprodLeftUnitor}}

\begin{defn}[Coproduct left unitor]
The isomorphism $0 + C \cong C$: nothing in the left summand.
\begin{itemize}
  \item Forward: $\mathrm{inl}\,e \mapsto e.\mathrm{elim}$ (impossible case);
    $\mathrm{inr}\,s \mapsto s$.
  \item Backward: $s \mapsto \mathrm{inr}\,s$.
\end{itemize}
\end{defn}

\paragraph{Lean.}
\begin{lstlisting}[style=lean]
def Container.coprodLeftUnitor (C : Container) :
    ContainerIso (Container.zero.coprod C) C where
  hom :=
    { onShapes := fun x => match x with | Sum.inl e => e.elim | Sum.inr s => s
      onPos    := fun x => match x with | Sum.inl e => e.elim | Sum.inr _ => id }
  inv :=
    { onShapes := Sum.inr
      onPos    := fun _ => id }
  hom_inv := ...;  inv_hom := ...
\end{lstlisting}

\paragraph{Haskell.}
\begin{lstlisting}[style=haskell]
coprodLeftUnitor :: ContainerIso (Coprod Zero c) c
-- fwd: Left (absurd) or Right s -> s
-- bwd: s -> Right s
\end{lstlisting}


\subsection{Definition: \lean{Container.coprodRightUnitor}}

\begin{defn}[Coproduct right unitor]
The isomorphism $C + 0 \cong C$: nothing in the right summand.
\end{defn}

\paragraph{Lean.}
\begin{lstlisting}[style=lean]
def Container.coprodRightUnitor (C : Container) :
    ContainerIso (C.coprod Container.zero) C where
  hom :=
    { onShapes := fun x => match x with | Sum.inl s => s | Sum.inr e => e.elim
      onPos    := fun x => match x with | Sum.inl _ => id | Sum.inr e => e.elim }
  inv := { onShapes := Sum.inl; onPos := fun _ => id }
  hom_inv := ...;  inv_hom := ...
\end{lstlisting}

\paragraph{Haskell.}
\begin{lstlisting}[style=haskell]
coprodRightUnitor :: ContainerIso (Coprod c Zero) c
\end{lstlisting}


\subsection{Theorem: \lean{Container.coprodAssociator\_naturality}}

\begin{thm}[Naturality of coproduct associator]
\[
  (\lean{coprod}_2\,(\lean{coprod}_2\,f_1\,f_2)\,f_3) \cdot \alpha^+_{Y_1,Y_2,Y_3}
  \;=\;
  \alpha^+_{X_1,X_2,X_3} \cdot \lean{coprod}_2\,f_1\,(\lean{coprod}_2\,f_2\,f_3).
\]
\end{thm}

\paragraph{Lean.}
\begin{lstlisting}[style=lean]
theorem Container.coprodAssociator_naturality
    {X1 Y1 X2 Y2 X3 Y3 : Container}
    (f1 : ContainerMorphism X1 Y1) (f2 : ContainerMorphism X2 Y2)
    (f3 : ContainerMorphism X3 Y3) :
    (Container.coprod2 (Container.coprod2 f1 f2) f3).comp
        (Container.coprodAssociator Y1 Y2 Y3).hom
      = (Container.coprodAssociator X1 X2 X3).hom.comp
          (Container.coprod2 f1 (Container.coprod2 f2 f3))
\end{lstlisting}

\paragraph{Haskell.}
\begin{lstlisting}[style=haskell]
-- coprodAssociator_naturality: naturality square for coprodAssociator
\end{lstlisting}


\subsection{Theorem: \lean{Container.coprodLeftUnitor\_naturality}}

\begin{thm}[Naturality of coproduct left unitor]
\[
  (\lean{coprod}_2\,(\mathrm{id}_0)\,f) \cdot \lambda^+_Y
  \;=\;
  \lambda^+_X \cdot f.
\]
\end{thm}

\paragraph{Lean.}
\begin{lstlisting}[style=lean]
theorem Container.coprodLeftUnitor_naturality {X Y : Container}
    (f : ContainerMorphism X Y) :
    (Container.coprod2 (ContainerMorphism.id Container.zero) f).comp
        (Container.coprodLeftUnitor Y).hom
      = (Container.coprodLeftUnitor X).hom.comp f
\end{lstlisting}


\subsection{Theorem: \lean{Container.coprodRightUnitor\_naturality}}

\begin{thm}[Naturality of coproduct right unitor]
\[
  (\lean{coprod}_2\,f\,(\mathrm{id}_0)) \cdot \rho^+_Y
  \;=\;
  \rho^+_X \cdot f.
\]
\end{thm}

\paragraph{Lean.}
\begin{lstlisting}[style=lean]
theorem Container.coprodRightUnitor_naturality {X Y : Container}
    (f : ContainerMorphism X Y) :
    (Container.coprod2 f (ContainerMorphism.id Container.zero)).comp
        (Container.coprodRightUnitor Y).hom
      = (Container.coprodRightUnitor X).hom.comp f
\end{lstlisting}


\subsection{Theorem: \lean{Container.coprodPentagon}}

\begin{thm}[Pentagon for $+$]
The two re-association paths of
$((C_1 + C_2) + C_3) + C_4 \to C_1 + (C_2 + (C_3 + C_4))$ agree.
\end{thm}

\paragraph{Lean.}
\begin{lstlisting}[style=lean]
theorem Container.coprodPentagon (C1 C2 C3 C4 : Container) :
    ((Container.coprod2 (Container.coprodAssociator C1 C2 C3).hom
          (ContainerMorphism.id C4)).comp
        (Container.coprodAssociator C1 (C2.coprod C3) C4).hom).comp
        (Container.coprod2 (ContainerMorphism.id C1)
          (Container.coprodAssociator C2 C3 C4).hom)
      = (Container.coprodAssociator (C1.coprod C2) C3 C4).hom.comp
          (Container.coprodAssociator C1 C2 (C3.coprod C4)).hom
\end{lstlisting}


\subsection{Theorem: \lean{Container.coprodTriangle}}

\begin{thm}[Triangle for $+$]
\[
  \alpha^+_{C_1,0,C_2} \cdot (\mathrm{id}_{C_1} + \lambda^+_{C_2})
  \;=\;
  \rho^+_{C_1} + \mathrm{id}_{C_2}.
\]
\end{thm}

\paragraph{Lean.}
\begin{lstlisting}[style=lean]
theorem Container.coprodTriangle (C1 C2 : Container) :
    (Container.coprodAssociator C1 Container.zero C2).hom.comp
        (Container.coprod2 (ContainerMorphism.id C1)
          (Container.coprodLeftUnitor C2).hom)
      = Container.coprod2 (Container.coprodRightUnitor C1).hom
          (ContainerMorphism.id C2)
\end{lstlisting}


\subsection{Definition: \lean{contCoprodMonoidal}}

\begin{defn}[$(\mathbf{Cont}, +, 0)$ is a monoidal category]
A named \lean{def} (not an \lean{instance}): the coproduct $+$ with morphism
action $\lean{coprod}_2$, unit $0$, and the structural isomorphisms above satisfy
all monoidal-category axioms.
\end{defn}

\paragraph{Lean.}
\begin{lstlisting}[style=lean]
@[reducible] def contCoprodMonoidal : MonoidalCategory Container where
  tensorObj         := Container.coprod
  tensorHom         := Container.coprod2
  tensorHom_id      := Container.coprod2_id
  tensorHom_comp    := fun f1 g1 f2 g2 => Container.coprod2_comp f1 g1 f2 g2
  tensorUnit        := Container.zero
  associator        := fun X Y Z => (Container.coprodAssociator X Y Z).toCatIso
  leftUnitor        := fun X => (Container.coprodLeftUnitor X).toCatIso
  rightUnitor       := fun X => (Container.coprodRightUnitor X).toCatIso
  associator_naturality  :=
    fun f1 f2 f3 => Container.coprodAssociator_naturality f1 f2 f3
  leftUnitor_naturality  := fun f => Container.coprodLeftUnitor_naturality f
  rightUnitor_naturality := fun f => Container.coprodRightUnitor_naturality f
  pentagon          := Container.coprodPentagon
  triangle          := Container.coprodTriangle
\end{lstlisting}

\paragraph{Haskell.}
\begin{lstlisting}[style=haskell]
contCoprodMonoidal :: MonoidalCategoryDef ContainerMorph
contCoprodMonoidal = MonoidalCategoryDef
  { tensor = coprod2, unit = zeroContainer, ... }
\end{lstlisting}


%===================================================================
\section{The categorical product \texorpdfstring{$(\times, 1)$}{(times, 1)}}
%===================================================================

\subsection{Definition: \lean{Container.prod\textsubscript{2}}}

\begin{defn}[Bifunctorial action of $\times$]
The morphism action $\lean{prod}_2\,\mu\,\nu : C \times D \Rightarrow C' \times D'$
is derived from the product's universal property: pair the compositions
$\pi_1 \cdot \mu$ and $\pi_2 \cdot \nu$.
\end{defn}

\paragraph{Lean.}
\begin{lstlisting}[style=lean]
def Container.prod2 {C C' D D' : Container}
    (mu : ContainerMorphism C C') (nu : ContainerMorphism D D') :
    ContainerMorphism (C.prod D) (C'.prod D') :=
  ((Container.fst C D).comp mu).pair ((Container.snd C D).comp nu)
\end{lstlisting}

\paragraph{Haskell.}
\begin{lstlisting}[style=haskell]
prod2 :: Morph c p c' p' -> Morph d q d' q'
      -> Morph (c `Prod` d) (Either p q) (c' `Prod` d') (Either p' q')
prod2 mu nu = pairMorph (fstMorph `compM` mu) (sndMorph `compM` nu)
\end{lstlisting}


\subsection{Theorem: \lean{Container.prod\textsubscript{2}\_id}}

\begin{thm}[Product bifunctoriality: identity]
$\lean{prod}_2\,(\mathrm{id}_C)\,(\mathrm{id}_D) = \mathrm{id}_{C \times D}$.
\end{thm}

\paragraph{Lean.}
\begin{lstlisting}[style=lean]
theorem Container.prod2_id (C D : Container) :
    Container.prod2 (ContainerMorphism.id C) (ContainerMorphism.id D)
      = ContainerMorphism.id (C.prod D)
\end{lstlisting}


\subsection{Theorem: \lean{Container.prod\textsubscript{2}\_comp}}

\begin{thm}[Product bifunctoriality: interchange]
$\lean{prod}_2\,(\mu \cdot \mu')\,(\nu \cdot \nu')
= \lean{prod}_2\,\mu\,\nu \;\cdot\; \lean{prod}_2\,\mu'\,\nu'$.
\end{thm}

\paragraph{Lean.}
\begin{lstlisting}[style=lean]
theorem Container.prod2_comp {C C' C'' D D' D'' : Container}
    (mu : ContainerMorphism C C') (mu' : ContainerMorphism C' C'')
    (nu : ContainerMorphism D D') (nu' : ContainerMorphism D' D'') :
    Container.prod2 (mu.comp mu') (nu.comp nu')
      = (Container.prod2 mu nu).comp (Container.prod2 mu' nu')
\end{lstlisting}


\subsection{Definition: \lean{Container.prodAssociator}}

\begin{defn}[Product associator]
The isomorphism $(C \times D) \times E \cong C \times (D \times E)$:
\begin{itemize}
  \item \emph{Shapes}: $\mathrm{Prod}$-reassociation
    $((s,t),u) \mapsto (s,(t,u))$.
  \item \emph{Positions}: $\mathrm{Sum}$-reassociation \emph{backwards}
    (since position maps are contravariant); the backward direction reassociates
    $P\,s \sqcup (Q\,t \sqcup R\,u)$ to $(P\,s \sqcup Q\,t) \sqcup R\,u$.
\end{itemize}
\end{defn}

\paragraph{Lean.}
\begin{lstlisting}[style=lean]
def Container.prodAssociator (C D E : Container) :
    ContainerIso ((C.prod D).prod E) (C.prod (D.prod E)) where
  hom :=
    { onShapes := fun x => (x.1.1, (x.1.2, x.2))
      onPos    := fun _ p =>
        Sum.elim (fun a => Sum.inl (Sum.inl a))
          (Sum.elim (fun b => Sum.inl (Sum.inr b))
                    (fun c => Sum.inr c)) p }
  inv :=
    { onShapes := fun x => ((x.1, x.2.1), x.2.2)
      onPos    := fun _ p =>
        Sum.elim (Sum.elim (fun a => Sum.inl a)
                           (fun b => Sum.inr (Sum.inl b)))
                 (fun c => Sum.inr (Sum.inr c)) p }
  hom_inv := ...;  inv_hom := ...
\end{lstlisting}

\paragraph{Haskell.}
\begin{lstlisting}[style=haskell]
-- Shapes: reassociate tuples.
-- Positions: reassociate Either (backwards, since contravariant).
prodAssociator :: ContainerIso (Prod (Prod c d) e) (Prod c (Prod d e))
\end{lstlisting}

\begin{fidelitynote}
The positions of $C \times D$ are $P\,s \sqcup Q\,t$.  The associator's position
map goes \emph{backward}: given a position of $C \times (D \times E)$, it
produces a position of $(C \times D) \times E$.  The backward position map of
the forward direction is a \emph{rightward} $\mathrm{Sum}$-reassociation
(injecting left-nested sums into right-nested ones), which is what the code
shows.
\end{fidelitynote}


\subsection{Definition: \lean{Container.prodLeftUnitor}}

\begin{defn}[Product left unitor]
The isomorphism $1 \times C \cong C$.  Since $1$ has shape $\mathbf{1}$ and no
positions, the position type of $1 \times C$ at shape $(\ast, s)$ is
$\emptyset \sqcup P\,s \cong P\,s$.
\begin{itemize}
  \item Forward on shapes: $(\ast, s) \mapsto s$.
  \item Forward on positions: inject right, $p \mapsto \mathrm{inr}\,p$.
\end{itemize}
\end{defn}

\paragraph{Lean.}
\begin{lstlisting}[style=lean]
def Container.prodLeftUnitor (C : Container) :
    ContainerIso (Container.one.prod C) C where
  hom :=
    { onShapes := fun x => x.2
      onPos    := fun _ p => Sum.inr p }
  inv :=
    { onShapes := fun s => ((), s)
      onPos    := fun _ p => Sum.elim (fun e => e.elim) id p }
  hom_inv := ...;  inv_hom := ...
\end{lstlisting}


\subsection{Definition: \lean{Container.prodRightUnitor}}

\begin{defn}[Product right unitor]
The isomorphism $C \times 1 \cong C$:
\begin{itemize}
  \item Forward on shapes: $(s, \ast) \mapsto s$.
  \item Forward on positions: inject left, $p \mapsto \mathrm{inl}\,p$.
\end{itemize}
\end{defn}

\paragraph{Lean.}
\begin{lstlisting}[style=lean]
def Container.prodRightUnitor (C : Container) :
    ContainerIso (C.prod Container.one) C where
  hom :=
    { onShapes := fun x => x.1
      onPos    := fun _ p => Sum.inl p }
  inv :=
    { onShapes := fun s => (s, ())
      onPos    := fun _ p => Sum.elim id (fun e => e.elim) p }
  hom_inv := ...;  inv_hom := ...
\end{lstlisting}


\subsection{Theorem: \lean{Container.prodAssociator\_naturality}}

\begin{thm}[Naturality of product associator]
\[
  (\lean{prod}_2\,(\lean{prod}_2\,f_1\,f_2)\,f_3) \cdot \alpha^\times_{Y_1,Y_2,Y_3}
  \;=\;
  \alpha^\times_{X_1,X_2,X_3} \cdot \lean{prod}_2\,f_1\,(\lean{prod}_2\,f_2\,f_3).
\]
\end{thm}

\paragraph{Lean.}
\begin{lstlisting}[style=lean]
theorem Container.prodAssociator_naturality
    {X1 Y1 X2 Y2 X3 Y3 : Container}
    (f1 : ContainerMorphism X1 Y1) (f2 : ContainerMorphism X2 Y2)
    (f3 : ContainerMorphism X3 Y3) :
    (Container.prod2 (Container.prod2 f1 f2) f3).comp
        (Container.prodAssociator Y1 Y2 Y3).hom
      = (Container.prodAssociator X1 X2 X3).hom.comp
          (Container.prod2 f1 (Container.prod2 f2 f3))
\end{lstlisting}


\subsection{Theorem: \lean{Container.prodLeftUnitor\_naturality}}

\begin{thm}[Naturality of product left unitor]
$(\lean{prod}_2\,(\mathrm{id}_1)\,f) \cdot \lambda^\times_Y = \lambda^\times_X \cdot f$.
\end{thm}

\paragraph{Lean.}
\begin{lstlisting}[style=lean]
theorem Container.prodLeftUnitor_naturality {X Y : Container}
    (f : ContainerMorphism X Y) :
    (Container.prod2 (ContainerMorphism.id Container.one) f).comp
        (Container.prodLeftUnitor Y).hom
      = (Container.prodLeftUnitor X).hom.comp f := rfl
\end{lstlisting}


\subsection{Theorem: \lean{Container.prodRightUnitor\_naturality}}

\begin{thm}[Naturality of product right unitor]
$(\lean{prod}_2\,f\,(\mathrm{id}_1)) \cdot \rho^\times_Y = \rho^\times_X \cdot f$.
\end{thm}

\paragraph{Lean.}
\begin{lstlisting}[style=lean]
theorem Container.prodRightUnitor_naturality {X Y : Container}
    (f : ContainerMorphism X Y) :
    (Container.prod2 f (ContainerMorphism.id Container.one)).comp
        (Container.prodRightUnitor Y).hom
      = (Container.prodRightUnitor X).hom.comp f := rfl
\end{lstlisting}


\subsection{Theorem: \lean{Container.prodPentagon}}

\begin{thm}[Pentagon for $\times$]
The two re-association paths of
$((C_1 \times C_2) \times C_3) \times C_4 \to C_1 \times (C_2 \times (C_3 \times C_4))$
agree.
\end{thm}

\paragraph{Lean.}
\begin{lstlisting}[style=lean]
theorem Container.prodPentagon (C1 C2 C3 C4 : Container) :
    ((Container.prod2 (Container.prodAssociator C1 C2 C3).hom
          (ContainerMorphism.id C4)).comp
        (Container.prodAssociator C1 (C2.prod C3) C4).hom).comp
        (Container.prod2 (ContainerMorphism.id C1)
          (Container.prodAssociator C2 C3 C4).hom)
      = (Container.prodAssociator (C1.prod C2) C3 C4).hom.comp
          (Container.prodAssociator C1 C2 (C3.prod C4)).hom
\end{lstlisting}


\subsection{Theorem: \lean{Container.prodTriangle}}

\begin{thm}[Triangle for $\times$]
$\alpha^\times_{C_1,1,C_2} \cdot (\mathrm{id}_{C_1} \times \lambda^\times_{C_2})
= \rho^\times_{C_1} \times \mathrm{id}_{C_2}$.
\end{thm}

\paragraph{Lean.}
\begin{lstlisting}[style=lean]
theorem Container.prodTriangle (C1 C2 : Container) :
    (Container.prodAssociator C1 Container.one C2).hom.comp
        (Container.prod2 (ContainerMorphism.id C1)
          (Container.prodLeftUnitor C2).hom)
      = Container.prod2 (Container.prodRightUnitor C1).hom
          (ContainerMorphism.id C2)
\end{lstlisting}


\subsection{Definition: \lean{contProdMonoidal}}

\begin{defn}[$(\mathbf{Cont}, \times, 1)$ is a monoidal category]
A named \lean{def}: the categorical product $\times$ with morphism action
$\lean{prod}_2$, unit $1$, and the structural isomorphisms satisfy all
monoidal-category axioms.
\end{defn}

\paragraph{Lean.}
\begin{lstlisting}[style=lean]
@[reducible] def contProdMonoidal : MonoidalCategory Container where
  tensorObj         := Container.prod
  tensorHom         := Container.prod2
  tensorHom_id      := Container.prod2_id
  tensorHom_comp    := fun f1 g1 f2 g2 => Container.prod2_comp f1 g1 f2 g2
  tensorUnit        := Container.one
  associator        := fun X Y Z => (Container.prodAssociator X Y Z).toCatIso
  leftUnitor        := fun X => (Container.prodLeftUnitor X).toCatIso
  rightUnitor       := fun X => (Container.prodRightUnitor X).toCatIso
  associator_naturality  :=
    fun f1 f2 f3 => Container.prodAssociator_naturality f1 f2 f3
  leftUnitor_naturality  := fun f => Container.prodLeftUnitor_naturality f
  rightUnitor_naturality := fun f => Container.prodRightUnitor_naturality f
  pentagon          := Container.prodPentagon
  triangle          := Container.prodTriangle
\end{lstlisting}


%===================================================================
\section{The extension functor as a strong monoidal functor}
%===================================================================

\subsection{Definitions and theorems: coproduct comparison}

\begin{defn}[Coproduct comparison isomorphism]
The maps
\[
  \lean{coprodExt} : \llbracket C + D \rrbracket X \xrightarrow{\sim}
    \llbracket C \rrbracket X \sqcup \llbracket D \rrbracket X,
\]
\[
  \lean{coprodExtInv} : \llbracket C \rrbracket X \sqcup \llbracket D \rrbracket X
    \xrightarrow{\sim} \llbracket C + D \rrbracket X.
\]
A $(C+D)$-structure is a $C$-structure or a $D$-structure, according to the
injection its shape came from.  Both round trips hold.
\end{defn}

\paragraph{Lean.}
\begin{lstlisting}[style=lean]
def Container.coprodExt {C D : Container} {X : Type} :
    Ext (C.coprod D) X -> Ext C X |+| Ext D X :=
  fun p => match p with
    | <|Sum.inl s, f|> => Sum.inl <|s, f|>
    | <|Sum.inr t, f|> => Sum.inr <|t, f|>

def Container.coprodExtInv {C D : Container} {X : Type} :
    Ext C X |+| Ext D X -> Ext (C.coprod D) X :=
  fun p => match p with
    | Sum.inl q => <|Sum.inl q.1, q.2|>
    | Sum.inr q => <|Sum.inr q.1, q.2|>

theorem Container.coprodExtInv_coprodExt {C D : Container} {X : Type}
    (p : Ext (C.coprod D) X) :
    Container.coprodExtInv (Container.coprodExt p) = p

theorem Container.coprodExt_coprodExtInv {C D : Container} {X : Type}
    (p : Ext C X |+| Ext D X) :
    Container.coprodExt (Container.coprodExtInv p) = p
\end{lstlisting}

\paragraph{Haskell.}
\begin{lstlisting}[style=haskell]
coprodExt :: Ext (Coprod c d) x -> Either (Ext c x) (Ext d x)
coprodExt (Ext (Left s)  f) = Left  (Ext s f)
coprodExt (Ext (Right t) f) = Right (Ext t f)

coprodExtInv :: Either (Ext c x) (Ext d x) -> Ext (Coprod c d) x
coprodExtInv (Left  (Ext s f)) = Ext (Left  s) f
coprodExtInv (Right (Ext t f)) = Ext (Right t) f
\end{lstlisting}


\subsection{Definition/Theorem: \lean{Container.zeroExt}}

\begin{defn}[Unit comparison for $0$]
The map $\llbracket 0 \rrbracket X \to \emptyset$: there is no shape, so there is
no structure.  The round trip is trivial (\lean{Empty.elim}).
\end{defn}

\paragraph{Lean.}
\begin{lstlisting}[style=lean]
def Container.zeroExt {X : Type} : Ext Container.zero X -> Empty :=
  fun p => p.1.elim

theorem Container.zeroExt_iso {X : Type} (p : Ext Container.zero X) :
    (Container.zeroExt p).elim = p := p.1.elim
\end{lstlisting}

\paragraph{Haskell.}
\begin{lstlisting}[style=haskell]
zeroExt :: Ext ZeroContainer x -> Void
zeroExt (Ext s _) = absurd s
\end{lstlisting}


\subsection{Definitions and theorems: product comparison}

\begin{defn}[Product comparison isomorphism]
The maps
\[
  \lean{prodExt} : \llbracket C \times D \rrbracket X \xrightarrow{\sim}
    \llbracket C \rrbracket X \times \llbracket D \rrbracket X,
\]
\[
  \lean{prodExtInv} : \llbracket C \rrbracket X \times \llbracket D \rrbracket X
    \xrightarrow{\sim} \llbracket C \times D \rrbracket X.
\]
Since positions of $C \times D$ are $P\,s \sqcup Q\,t$, a labelling
$P\,s \sqcup Q\,t \to X$ is a pair of labellings $(P\,s \to X, Q\,t \to X)$
by the exponential law $y^{A+B} \cong y^A \times y^B$.  This is why the
product container has the \emph{pointwise product} of extensions.
\end{defn}

\paragraph{Lean.}
\begin{lstlisting}[style=lean]
def Container.prodExt {C D : Container} {X : Type} :
    Ext (C.prod D) X -> Ext C X * Ext D X :=
  fun p => (<|p.1.1, fun i => p.2 (Sum.inl i)|>,
            <|p.1.2, fun j => p.2 (Sum.inr j)|>)

def Container.prodExtInv {C D : Container} {X : Type} :
    Ext C X * Ext D X -> Ext (C.prod D) X :=
  fun p => <|(p.1.1, p.2.1), Sum.elim p.1.2 p.2.2|>

theorem Container.prodExtInv_prodExt {C D : Container} {X : Type}
    (p : Ext (C.prod D) X) :
    Container.prodExtInv (Container.prodExt p) = p

theorem Container.prodExt_prodExtInv {C D : Container} {X : Type}
    (p : Ext C X * Ext D X) :
    Container.prodExt (Container.prodExtInv p) = p
\end{lstlisting}

\paragraph{Haskell.}
\begin{lstlisting}[style=haskell]
prodExt :: Ext (Prod c d) x -> (Ext c x, Ext d x)
prodExt (Ext (s, t) f) =
  (Ext s (\i -> f (Left i)), Ext t (\j -> f (Right j)))

prodExtInv :: (Ext c x, Ext d x) -> Ext (Prod c d) x
prodExtInv (Ext s f, Ext t g) = Ext (s, t) (either f g)
\end{lstlisting}


\subsection{Definitions and theorems: unit comparison for $1$}

\begin{defn}[Unit comparison for $1$]
The isomorphism $\llbracket 1 \rrbracket X \cong \mathbf{1}$:
one shape, no positions, so exactly one structure.
\end{defn}

\paragraph{Lean.}
\begin{lstlisting}[style=lean]
def Container.oneExt {X : Type} : Ext Container.one X -> Unit :=
  fun _ => ()

def Container.oneExtInv {X : Type} : Unit -> Ext Container.one X :=
  fun _ => <|(), fun e => e.elim|>

theorem Container.oneExtInv_oneExt {X : Type}
    (p : Ext Container.one X) :
    Container.oneExtInv (Container.oneExt p) = p

theorem Container.oneExt_oneExtInv {X : Type} (u : Unit) :
    Container.oneExt (Container.oneExtInv (X := X) u) = u := rfl
\end{lstlisting}

\paragraph{Haskell.}
\begin{lstlisting}[style=haskell]
oneExt    :: Ext OneContainer x -> ()
oneExt _  = ()

oneExtInv :: () -> Ext OneContainer x
oneExtInv () = Ext () absurd
\end{lstlisting}


\subsection{Definitions: shared unit comparison \texorpdfstring{$\llbracket y \rrbracket X \cong X$}{[y] X = X}}

\begin{defn}[Unit comparison for $y = I$]
The isomorphism $\llbracket y \rrbracket X \cong X$: read the label at the unique
position.  This is the unit comparison for both $\otimes$ and $\triangleleft$,
since $y = I$.
\[
  \lean{yExt}(s, v) = v(\ast), \qquad
  \lean{yExtInv}(x) = (\ast,\; \_\mapsto x).
\]
\end{defn}

\paragraph{Lean.}
\begin{lstlisting}[style=lean]
def Container.yExt {X : Type} : Ext Container.y X -> X :=
  fun p => p.2 ()

def Container.yExtInv {X : Type} : X -> Ext Container.y X :=
  fun x => <|(), fun _ => x|>

theorem Container.yExtInv_yExt {X : Type}
    (p : Ext Container.y X) :
    Container.yExtInv (Container.yExt p) = p := rfl

theorem Container.yExt_yExtInv {X : Type} (x : X) :
    Container.yExt (Container.yExtInv x) = x := rfl
\end{lstlisting}

\paragraph{Haskell.}
\begin{lstlisting}[style=haskell]
yExt :: Ext YContainer x -> x
yExt (Ext () v) = v ()

yExtInv :: x -> Ext YContainer x
yExtInv x = Ext () (const x)
\end{lstlisting}

\begin{fidelitynote}
Both round trips hold by \lean{rfl}: definitional $\eta$ for \lean{Unit} and the
\lean{Sigma} / \lean{Ext} type.
\end{fidelitynote}

\end{document}
